\section{文献阅读}
\label{sec:literature-review}

除程序实现和算例分析外，本项目还围绕潮流计算、病态机理与静态电压稳定阅读了相关
书籍和论文。以下简要介绍对本文影响较大的内容。

\subsection{《高等电力网络分析》（张伯明、严正）}

该书帮助我系统梳理了潮流计算的数学模型、基本算法及发展过程，并补充了特殊潮流、
随机潮流和潮流跟踪等扩展内容。其中，连续潮流法的详细介绍打下基础。

另一个很有趣的内容是快速解耦法中的迭代取用顺序将内涵标准NR法迭代的主要耦合影响因素，
为快速解耦法的广泛适用范围做出解释。

\subsection{《雅可比矩阵对直角坐标牛顿法潮流计算收敛性影响的研究》}

该文的背景介绍部分对病态潮流的发展和相关研究进行了较完整的梳理，
提供了零散的病态原因以及病态算法，
本文进一步归纳出病态成因的物理可解性、初值吸引域和数值迭代过程的三级分类。

\subsection{《基于张量法的电力系统病态潮流求解与 \(P\)--\(V\) 曲线》}

该文介绍了不同于 NR 和快速解耦法的张量潮流方法。其对非线性模型和迭代格式的扩展
幅度大于最优乘子法，为理解病态潮流中改善收敛性的其他技术路线提供了对照。

\subsection{《基于连续潮流的电力系统电压稳定性研究》}

该文与《高等电力网络分析》共同加深了本文对标准连续潮流流程的认识，并重点启发了
对自然、弧长和绝对角度等不同参数化方式的比较，本文进一步适用了切线角度进行了尝试分析。

\subsection{《电力系统潮流的病态性及其算法研究》}

该文从同伦方法的角度解释了连续潮流的数学基础，使本文认识到 CPF 可以看作同伦思想
在潮流计算中的一种具体形式；通过调整同伦函数和延拓参数，还可以进一步改善复杂工况
下的跟踪与收敛性能。
